\setcounter{section}{12}
\section{Szereg Fouriera: współczynniki, ortogonalność i zbieżność}
\label{sec:dt-series-fourier}
\providecommand{\atlasref}[2]{#2}

Por. \atlasref{PDE-06::thm:pde6-fejer}{RRC-06: średnie Fejéra i szeregi Fouriera}.


Przez $\mathbb T=\R/(2\pi\mathbb Z)$ rozumiemy okrąg z
współrzędną $x\in[-\pi,\pi]$ i utożsamionymi końcami.
Funkcję z $L^1(-\pi,\pi)$ utożsamiamy z jej
$2\pi$-okresowym przedłużeniem, określonym z dokładnością
do zbioru miary zero; wartości na końcach przedziału
nie wpływają na całki. Dla takiej funkcji definiujemy
$\widehat f(k)=(2\pi)^{-1}\int_{-\pi}^{\pi}f(x)e^{-ikx}\,\mathrm d\lambda_1(x)$,
$k\in\mathbb Z$, oraz sumę częściową
$S_Nf(x)=\sum_{|k|\le N}\widehat f(k)e^{ikx}$.
W tym rozdziale $L^2(\mathbb T)$ oznacza zespoloną przestrzeń Hilberta
z iloczynem skalarnym
$\langle f,g\rangle=(2\pi)^{-1}\int_{-\pi}^{\pi}f(x)\overline{g(x)}\,\mathrm d\lambda_1(x)$.
Pojęcie całki Lebesgue'a i podstawowe własności $L^p$
traktujemy jako znane. Na okręgu używamy miary Haara
unormowanej do masy $1$, czyli $\mathrm d\lambda_1/(2\pi)$.
Na torusie $\mathbb T^d$ jej odpowiednikiem jest
$\mathrm d\lambda_d/(2\pi)^d$ na kostce fundamentalnej;
współczynniki Fouriera zawierają zatem czynnik $(2\pi)^{-d}$.
Norma $L^2$ względem nieznormalizowanej miary jest
$(2\pi)^{d/2}$ razy większa. Nie należy utożsamiać tej
normalizacji okresowej z unitarną normalizacją całki
Fouriera na $\R^d$ z definicji \ref{def:fourier-schwartz17}.
Ponieważ miara unormowana na
$\mathbb T$ ma masę $1$, nierówność Cauchy'ego--Schwarza
daje $(2\pi)^{-1}\|f\|_{L^1(-\pi,\pi)}\leq\|f\|_2$,
więc współczynniki Fouriera są określone dla każdego
$f\in L^2(\mathbb T)$.

\begin{lemat}[Ortogonalność i nierówność Bessela]
\label{lem:dt-series-bessel}
Niech $f\in L^2(\mathbb T)$ oraz $N\geq0$ będzie całkowite. Wówczas
$\langle e^{ikx},e^{i\ell x}\rangle=\delta_{k\ell}$,
$\|S_Nf\|_2^2=\sum_{|k|\le N}|\widehat f(k)|^2$ i
$\sum_{|k|\le N}|\widehat f(k)|^2\le\|f\|_2^2$.
\end{lemat}
\begin{proof}
Dla $k=\ell$ całka z jedynki względem miary
unormowanej wynosi $1$. Jeśli $k\ne\ell$, to
\[
 \frac1{2\pi}\int_{-\pi}^{\pi}e^{i(k-\ell)x}\,
 \mathrm d\lambda_1(x)
 =\frac{e^{i(k-\ell)\pi}-e^{-i(k-\ell)\pi}}
 {2\pi i(k-\ell)}=0,
\]
ponieważ $k-\ell$ jest całkowite. Otrzymujemy
ortonormalność. Rozwijając kwadrat normy skończonej
sumy $S_Nf$ i korzystając z tej ortonormalności,
dostajemy $\|S_Nf\|_2^2=\sum_{|k|\le N}|\widehat f(k)|^2$.
Ponadto dla $|\ell|\leq N$
\[
 \langle f-S_Nf,e^{i\ell x}\rangle
 =\widehat f(\ell)-\sum_{|k|\leq N}
 \widehat f(k)\delta_{k\ell}=0.
\]
Zatem $S_N$ jest rzutem ortogonalnym na przestrzeń
rozpiętą przez te wykładniki. W szczególności
$\langle f-S_Nf,S_Nf\rangle=0$, skąd
\[
 \|f\|_2^2=\|S_Nf\|_2^2+\|f-S_Nf\|_2^2.
\]
Nieujemność ostatniego składnika daje nierówność Bessela.
\end{proof}

Nie można jednak wywnioskować z ortogonalności, że $S_Nf\to f$.
Do tego potrzebujemy zupełności układu wykładniczego.

\begin{lemat}[Jądro Fejéra]
\label{lem:dt-series-fejer}
Dla całkowitego $N\geq0$ niech
\begin{equation}\label{eq:dt-series-fejer}
 F_N(t)=\frac1{N+1}\left|\sum_{j=0}^Ne^{ijt}\right|^2
 =\frac1{N+1}\left(\frac{\sin((N+1)t/2)}{\sin(t/2)}\right)^2
\end{equation}
(w punktach $t\in2\pi\mathbb Z$ przyjmujemy
$F_N(t)=N+1$, zgodnie z pierwszą postacią wzoru). Wówczas $F_N\ge0$,
$(2\pi)^{-1}\int_{-\pi}^{\pi}F_N(t)\,\mathrm d\lambda_1(t)=1$ oraz
dla każdego $\delta\in(0,\pi)$
$\int_{\delta\le|t|\le\pi}F_N(t)\,\mathrm d\lambda_1(t)\to0$.
\end{lemat}
\begin{proof}
Nieujemność jest bezpośrednia z definicji.
Dla $t\notin2\pi\mathbb Z$ z sumy geometrycznej
$\sum_{j=0}^N e^{ijt}=(1-e^{i(N+1)t})/(1-e^{it})$
i tożsamości $|1-e^{i\alpha}|=2|\sin(\alpha/2)|$
otrzymujemy drugą postać~\eqref{eq:dt-series-fejer}.
Rozwijając kwadrat modułu, mamy
\[
 F_N(t)=\frac1{N+1}\sum_{j,\ell=0}^N e^{i(j-\ell)t}
 =\sum_{|k|\leq N}\left(1-\frac{|k|}{N+1}\right)e^{ikt}.
\]
Dla ustalonego $|k|\leq N$ istnieje dokładnie
$N+1-|k|$ par $(j,\ell)$ z różnicą $j-\ell=k$:
dla $k\geq0$ wybieramy $\ell=0,\ldots,N-k$,
a dla $k<0$ zamieniamy role indeksów.
Wyraz stały jest równy $1$, a pozostałe
wykładniki mają całkę zero na pełnym okresie,
skąd normalizacja. Gdy $\delta\leq|t|\leq\pi$,
mamy $|\sin(t/2)|\geq\sin(\delta/2)>0$.
Dlatego
\[
 \int_{\delta\leq|t|\leq\pi}F_N(t)\,\mathrm d\lambda_1(t)
 \leq\frac{2\pi}{(N+1)\sin^2(\delta/2)}
 \longrightarrow0.
\]
\end{proof}

\begin{twierdzenie}[Fejér i zupełność układu trygonometrycznego]
\label{thm:dt-series-fejer}
Niech $f\in C(\mathbb T)$. Średnie Cesàra
$\sigma_N f=(N+1)^{-1}\sum_{j=0}^N S_jf$
zbieżne są jednostajnie do $f$ na $\mathbb T$.
W konsekwencji wielomiany trygonometryczne są gęste
w $C(\mathbb T)$ w normie supremum i w $L^2(\mathbb T)$
w normie $L^2$.
\end{twierdzenie}
\begin{proof}
W średniej z sum częściowych współczynnik
przy $\widehat f(k)e^{ikx}$ wynosi
$(N+1-|k|)/(N+1)$ dla $|k|\leq N$, ponieważ
wyraz ten występuje dla $j=|k|,\ldots,N$.
Z drugiej strony podstawienie $y=x-t$ daje
\[
 \frac1{2\pi}\int_{-\pi}^{\pi}e^{ikt}f(x-t)\,
 \mathrm d\lambda_1(t)=e^{ikx}\widehat f(k).
\]
Po podstawieniu całkujemy po odcinku
$[x-\pi,x+\pi]$, który możemy zastąpić
$[-\pi,\pi]$, gdyż $f(y)e^{-iky}$ jest
$2\pi$-okresowa i całka po pełnym okresie
nie zależy od jego początku (przez podział
odcinka i przesunięcie o wielokrotność okresu).
Rozwinięcie skończone z lematu~\ref{lem:dt-series-fejer}
daje więc
\[
 \sigma_Nf(x)=\frac1{2\pi}\int_{-\pi}^{\pi}
 F_N(t)f(x-t)\,\mathrm d\lambda_1(t).
\]
Odejmujemy $f(x)$ i dzielimy całkę na $|t|<\delta$
oraz $\delta\le|t|\le\pi$. Z jednostajnej ciągłości $f$
możemy dobrać $\delta$ tak, aby pierwszy składnik
miał moduł co najwyżej $\varepsilon$ dla wszystkich $x,N$.
Drugi ma moduł nie większy niż
$2\|f\|_\infty(2\pi)^{-1}\int_{|t|\ge\delta}F_N(t)
\,\mathrm d\lambda_1(t)$, co dąży do zera.
Ponieważ $F_N$ ma unormowaną całkę $1$, pierwszy
składnik jest rzeczywiście ograniczony przez
$\varepsilon$, niezależnie od $N$. Zatem dla
dostatecznie dużych $N$ norma supremum różnicy
nie przekracza $2\varepsilon$. Dowodzi to
jednostajnej zbieżności.

Uzasadnijmy również gęstość w $L^2$ z uwzględnieniem
okresowości. Funkcje proste są gęste w $L^2$
na przedziale skończonej miary. Dla mierzalnego
$A\subset(-\pi,\pi)$ oraz $\eta>0$ regularność
miary Lebesgue'a daje zbiór zwarty $K\subset A$
i otwarty $O\supset A$, $O\subset(-\pi,\pi)$,
takie że $\lambda_1(O\setminus K)<\eta$.
Gdy $K$ jest niepusty, wybieramy
$0<r<\operatorname{dist}(K,\R\setminus O)$.
Funkcja
\[
 \varphi(x)=\max\{0,1-\operatorname{dist}(x,K)/r\}
\]
jest ciągła, równa $1$ na $K$, ma zwarty nośnik
zawarty w $O$ i wartości w $[0,1]$.
Jeśli $K=\varnothing$, bierzemy $\varphi=0$.
W obu przypadkach
$\|\chi_A-\varphi\|_2^2\leq\eta/(2\pi)$.
Przedłużenie $\varphi$ okresowo jest ciągłe,
bo $\varphi$ znika w otoczeniu końców przedziału.
Aproksymując w ten sposób każdy składnik
funkcji prostej i stosując nierówność trójkąta,
otrzymujemy gęstość $C(\mathbb T)$ w $L^2(\mathbb T)$.
Dla $f\in L^2$ i $\varepsilon>0$ wybieramy
$g\in C(\mathbb T)$ z $\|f-g\|_2<\varepsilon/2$,
a następnie wielomian trygonometryczny $p$
z $\|g-p\|_\infty<\varepsilon/2$.
Ponieważ $\|g-p\|_2\leq\|g-p\|_\infty$,
mamy $\|f-p\|_2<\varepsilon$.
\end{proof}

\begin{twierdzenie}[Zbieżność w $L^2$ i tożsamość Parsevala]
\label{thm:dt-series-parseval}
Dla dowolnego $f\in L^2(\mathbb T)$ zachodzą
$S_Nf\to f$ w normie $L^2$ oraz
\begin{equation}\label{eq:dt-series-parseval}
 \frac1{2\pi}\int_{-\pi}^{\pi}|f(x)|^2\,\mathrm d\lambda_1(x)
 =\sum_{k\in\mathbb Z}|\widehat f(k)|^2.
\end{equation}
\end{twierdzenie}
\begin{proof}
Dla $\varepsilon>0$ wybieramy wielomian trygonometryczny $p$
z $\|f-p\|_2<\varepsilon$. Dla $N$ nie mniejszych niż stopień
$p$ mamy $S_Np=p$. Ponieważ $f-S_Nf$ jest
prostopadła do przestrzeni, do której należy
$S_Nf-p$, zachodzi
\[
 \|f-p\|_2^2=\|f-S_Nf\|_2^2+\|S_Nf-p\|_2^2.
\]
Zatem $\|f-S_Nf\|_2\leq\|f-p\|_2<\varepsilon$,
co dowodzi zbieżności.
W konsekwencji $\|S_Nf\|_2\to\|f\|_2$,
a z tożsamości $\|S_Nf\|_2^2=
\sum_{|k|\leq N}|\widehat f(k)|^2$
w lemacie~\ref{lem:dt-series-bessel} otrzymujemy
\eqref{eq:dt-series-parseval}, ponieważ sumy
nieujemnych wyrazów rosną do pełnej sumy po $\mathbb Z$.
\end{proof}

\begin{uwaga}[Jakiego rodzaju zbieżność została udowodniona?]
Sama zbieżność w normie $L^2$ nie pociąga
za sobą zbieżności jednostajnej ani punktowej
całego ciągu. Nie rozstrzygamy tutaj odrębnego
zagadnienia zbieżności prawie wszędzie sum
częściowych Fouriera. Dla funkcji ciągłych jednostajnie zbieżne są
\emph{średnie Fejéra}, niekoniecznie zwykłe sumy $S_N$.
Dla gładkich funkcji zachodzi dodatkowe szybkie zanikanie
współczynników, co wykażemy poniżej.
\end{uwaga}

\begin{lemat}[Zanikanie współczynników gładkiej funkcji]
\label{lem:dt-series-decay}
Jeśli $f\in C^m(\mathbb T)$, $m\geq1$ całkowite, to dla $k\ne0$
$|\widehat f(k)|\le\|f^{(m)}\|_{L^1(-\pi,\pi)}/(2\pi|k|^m)$.
Jeśli $f\in C^\infty(\mathbb T)$, to szereg Fouriera wraz
ze wszystkimi szeregami otrzymanymi przez różniczkowanie
zbiega jednostajnie.
\end{lemat}
\begin{proof}
Dla $k\ne0$ całkowanie przez części daje
\[
 \widehat f(k)=\frac1{ik}\widehat{f'}(k),
\]
ponieważ $f(\pi)e^{-ik\pi}=f(-\pi)e^{ik\pi}$.
Powtarzając ten rachunek $m$ razy (pochodne
do rzędu $m-1$ mają zgodne wartości na końcach
przedziału), dostajemy
$\widehat f(k)=(ik)^{-m}\widehat{f^{(m)}}(k)$,
co po oszacowaniu modułu całki daje żądaną nierówność.

Dla $f\in C^\infty$ i ustalonego całkowitego
$r\geq0$ wybieramy $m\geq r+2$. Wówczas
\[
 \sum_{k\ne0}|k|^r|\widehat f(k)|
 \leq\frac{\|f^{(m)}\|_{L^1(-\pi,\pi)}}{2\pi}
 \sum_{k\ne0}|k|^{r-m}<\infty.
\]
Kryterium Weierstrassa daje bezwzględną
i jednostajną zbieżność szeregu
$\sum_{k\in\mathbb Z}(ik)^r\widehat f(k)e^{ikx}$,
przy czym dla $r=0$ wyraz $k=0$ oznacza
$\widehat f(0)$, a dla $r>0$ wynosi zero.
Oznaczmy jego ciągłą sumę przez $G_r$.
Z całkowania przez części także dla $k=0$
(jeśli $r>0$, całka z $f^{(r)}$ po okresie
jest zerowa) współczynniki tego szeregu
są dokładnie współczynnikami Fouriera $f^{(r)}$.
Twierdzenie~\ref{thm:dt-series-parseval} daje więc
zbieżność tych samych sum częściowych do $f^{(r)}$
w normie $L^2$. Z jednostajnej zbieżności
mamy zarazem zbieżność w $L^2$ do $G_r$.
Jednoznaczność granicy w $L^2$ daje
$G_r=f^{(r)}$ prawie wszędzie, a ponieważ
obie funkcje są ciągłe, równość zachodzi
w każdym punkcie: niezerowa różnica
w jednym punkcie byłaby niezerowa na
pewnym otwartym łuku dodatniej miary.
W szczególności $G_0=f$ i wszystkie
różniczkowania wyraz po wyrazie prowadzą
do właściwych pochodnych.
\end{proof}
